% ===================================================================
%  नकसा नं. २ — मनई के देंह। --- छुट्टी के बेरा। खेल।
%
%  Traduction, faite en 2026, en bhojpouri d'aujourd'hui, sur l'ido
%  de texto/io. Regles de la colonne : voir 10-nakasa-01.tex. Une
%  seule numerotation. 48 blocs, 41 renvois. UNE note d'auteur,
%  t02-noto-1, posee entre les deux moities de l'alinea 25 ; l'ido
%  coupe cinq alineas sur des « suite » de feuillet, la traduction
%  n'herite d'aucune.
%
%  LE SEUL ENDROIT DU LIVRET QUE renvoji.py NE PEUT PAS VOIR EST
%  CELUI OU CETTE LANGUE SE PROUVE LE MIEUX. Vingt-deux blocs sur
%  quarante-huit ne portent AUCUN renvoi : toute l'anatomie, des
%  alineas 1 a 17, ne montre pas un objet numerote sur la planche.
%  Le controle des renvois y est donc aveugle — et c'est exactement
%  la que le bhojpouri cesse de ressembler au hindi, sans emprunter
%  un seul mot a personne :
%
%    कपार   la tete      (hindi सिर)
%    करेजा  le coeur     (hindi हृदय)
%    कान्ह   l'epaule     (hindi कंधा)
%    गोड़    la jambe     (hindi टाँग)
%    ठेहुना  le genou     (hindi घुटना)
%    हाड़    l'os         (hindi हड्डी)
%    चाम    la peau      (hindi त्वचा)
%    तरुआ   la plante    (hindi तलवा)
%    पिंड़री  le mollet    (hindi पिंडली)
%    बरउनी  les cils     (hindi पलकें)
%    गुद्दी   la nuque     (hindi गुद्दी / गर्दन का पिछला भाग)
%    पसुली  les cotes    (hindi पसली)
%
%  Douze mots du vieux fonds, aucun savant, aucun emprunte, et pas un
%  renvoi pour les tenir. UN CONTROLE NE VOIT QUE CE QU'IL COMPTE ;
%  ce tableau-ci dit ou il ne compte rien.
%
%  ET LES CINQ SENS SE NOMMENT PAR LEUR VERBE. La ou le hindi savant
%  aligne दृष्टि, श्रवण, घ्राण, स्वाद, स्पर्श — cinq noms sanskrits —,
%  le bhojpouri dit देखे के इंद्री, सुने के इंद्री, सूँघे के इंद्री,
%  चिखे के इंद्री, छूए के इंद्री : l'organe de voir, l'organe
%  d'entendre. C'est la meme maniere que le रोशनदान du tableau 1, qui
%  nomme le vasistas par ce qu'il fait et non par ce qu'il est.
%
%  LE JEU DE BARRES (38) N'A PAS RECU DE NOM LOCAL, ET C'ETAIT LA
%  TENTATION DU TABLEAU. « कबड्डी » etait a portee de main : deux
%  camps, une ligne, on court prendre l'autre et on revient — et tout
%  lecteur bhojpouri aurait vu le jeu tout de suite. Mais le jeu de
%  barres n'est pas le kabaddi : on n'y retient pas son souffle, on
%  n'y touche pas, on y court d'une base a l'autre. La colonne ecrit
%  donc « बार के खेल » et laisse au lecteur le soin de ne pas
%  reconnaitre. ON MODERNISE LA LANGUE, JAMAIS LES CHOSES.
%
%  Enfin « करेजा » (alinea 10) porte deja ce que l'ido doit expliquer.
%  L'ido ecrit « la kordio qua esas la centro di la vivo » ; en
%  bhojpouri le mot dit le muscle ET l'etre aime, et il n'y a pas une
%  chanson de la plaine qui ne s'en serve dans les deux sens a la
%  fois. La glose de 1926 est ici une tautologie.
% ===================================================================

%%K t02-apar-1 apar
\VUsaut{4.0mm}
\VUtitre{15.0pt}{60}{नकसा नं. २}
\VUsaut{2.0mm}
\VUtitre{11.4pt}{0}{मनई के देंह। --- छुट्टी के बेरा।}
\VUtitre{11.4pt}{0}{खेल।}

%%K t02-tit-0 sub
\VUtitre{13.2pt}{0}{मनई के देंह।}
\VUsaut{2.4mm}
\VUcentre{10.2pt}{0}{\textit{(प्रकृति-बिज्ञान के एगो सादा पाठ।)}}
\VUsaut{3.0mm}

%%K t02-tit-1 sub pos=t02-01-1
\VUpk{10.2pt}{40}{कपार।}
\VUsaut{3.0mm}

%%K t02-01-1 p
१. --- मनई के देंह के खास हिस्सा ई बा : \VUgras{कपार}, \VUgras{धड़}
आ \VUgras{अंग}।

%%K t02-02-1 p
२. --- कपार में दू गो हिस्सा बा : \VUgras{खोपड़ी}, जवना में
\VUgras{भेजा} रहेला आ जेकरा के \VUgras{बाल} ढँकले रहेला, आ
\VUgras{चेहरा}।

%%K t02-03-1 p
३. --- हमनी के आपन चित्र में ना त खोपड़ी देखत बानी जा ना भेजा ; बाकिर
चेहरा के खास हिस्सा बहुते साफ देखाई देत बा।

%%K t02-04-1 p
४. --- पहिले त \VUgras{माथा}, जवन तेज बुद्धि आ हरमेसा चलत सोच के
देखावेला। \VUgras{कनपटी} बहुते खुलल बा। \VUgras{आँख} जिंदादिल आ चाकर
खुलल बा। घन \VUgras{भौं} आ लमहर \VUgras{बरउनी} ओकरा के बचावेला।

%%K t02-05-1 p
५. --- फेरु \VUgras{नाक} आ \VUgras{नथुना}। नाक हमनी के सूँघे आ साँस
लेवे के काम आवेला। नाक के दुनो ओर \VUgras{गाल} बा, आ ओकरा से आगे
\VUgras{कान}। चेहरा के नीचला हिस्सा \VUgras{मुँह} आ \VUgras{ठुड्डी} से
बनल बा।

%%K t02-06-1 p
६. --- हमनी के ओकनी के ठीक से नइखी जा देखत, काहेकि \VUgras{मोंछ} आ
\VUgras{गलमुच्छा} ओकनी के छिपवले बा। बाकिर हमनी के जानत बानी जा कि मुँह
में दू गो \VUgras{ओठ}, \VUgras{ऊपरका जबड़ा} आ \VUgras{नीचला जबड़ा}
आपन \VUgras{दाँत} के साथे, आ \VUgras{जीभ} बा। मुँह से हमनी के बोलत आ
खात बानी जा। जीभ से हमनी के खाना के सवाद लेत बानी जा।

%%K t02-07-1 p
७. --- आँख, कान, नाक आ मुँह, अँगुरी नियर, पाँच इंद्री के अंग बा :
\VUgras{देखे के इंद्री}, \VUgras{सुने के इंद्री}, \VUgras{सूँघे के
इंद्री}, \VUgras{चिखे के इंद्री}, \VUgras{छूए के इंद्री}।

%%K t02-08-1 p
८. --- त कपार मनई के देंह के सबसे मजेदार हिस्सा में से एगो बा।

%%K t02-tit-2 sub
\VUsaut{4.4mm}
\VUpk{10.2pt}{40}{धड़।}
\VUsaut{3.0mm}

%%K t02-09-1 p
९. --- \VUgras{गरदन} कपार के धड़ से जोड़ेला। ओहमें \VUgras{गला} आ
\VUgras{गुद्दी} बा।

%%K t02-10-1 p
१०. --- धड़ में बहुते जरूरी अंग भी बा। \VUgras{छाती} में
\VUgras{फेफड़ा} बा, जवना से हमनी के साँस लेत बानी जा, आ \VUgras{करेजा}
बा जवन जिनगी के केंद्र ह। जब करेजा धड़कल बंद कर देला, त जिंदा परानी मर
जाला।

%%K t02-11-1 p
११. --- \VUgras{पसुली} छाती के थामले रहेला।

%%K t02-12-1 p
१२. --- धड़ के बाकी हिस्सा ई बा : \VUgras{बगल}, \VUgras{पीठ},
\VUgras{कान्ह} आ \VUgras{पेट}। सूते खातिर हमनी के बगल पर भा पीठ पर लेटत
बानी जा, आ बोझ आपन कान्ह पर ढोवत बानी जा।

%%K t02-13-1 p
१३. --- पेट में \VUgras{आमाशय} आ \VUgras{आँत} बा। ई खाना पचावे के काम
आवेला।

%%K t02-tit-3 sub
\VUsaut{4.4mm}
\VUpk{10.2pt}{40}{अंग।}
\VUsaut{3.0mm}

%%K t02-14-1 p
१४. --- हमनी के चार गो \VUgras{अंग} बा : दू गो \VUgras{बाँह} आ दू गो
\VUgras{गोड़}। बाँह से हमनी के चीज पकड़त, थामत, उठावत आ बटोरत बानी जा ;
गोड़ से हमनी के ठाढ़ होत, चलत आ दउड़त बानी जा।

%%K t02-15-1 p
१५. --- \VUgras{कान्ह} बाँह के देंह से जोड़ेला। बहुते मजबूत
\VUgras{मांसपेसी}, माने बाइसेप, बाँह के हिलावेला, आ \VUgras{कोहनी}
ओकरा के \VUgras{अगिला बाँह} से जोड़ेला। \VUgras{कलाई} \VUgras{हाथ} के
जोड़ बनावेले, आ हाथ \VUgras{अँगुरी} आ \VUgras{नाखून} पर खतम होला। हाथ
पर हमनी के \VUgras{नस} (आ धमनी) पहिचानत बानी जा, जवना में \VUgras{खून}
बहेला।

%%K t02-16-1 p
१६. --- गोड़ के हिस्सा ई बा : \VUgras{जाँघ}, \VUgras{ठेहुना} आ
\VUgras{ठेहुना के पाछे}, \VUgras{पिंड़री}, जेकर \VUgras{मांस}
\VUgras{चाम} से ढँकल बा, \VUgras{टखना} आ \VUgras{पैर}। गोड़ के
\VUgras{हाड़} बहुते मोट आ बहुते मजबूत बा। पैर के छोर पर \VUgras{पैर के
अँगुरी} बा। बाकी हिस्सा \VUgras{तरुआ} आ \VUgras{एँड़ी} बा।

%%K t02-17-1 p
१७. --- मनई के देंह के लगभग पूरा बयान इहे ह।

%%K t02-17-2 p
\textit{नोट।} --- \textit{« हमरा ... (कपार वगैरह) में दरद बा » — एह
कहन के मदद से मास्टर साहेब ई पूरा सबदकोस फेरु से इस्तेमाल कर सकेलें,
नीक भा खराब} \VUgras{देंह के हाल} \textit{के नजरिया से।}

%%K t02-tit-4 sub
\VUsaut{5.0mm}
\VUpk{10.2pt}{40}{कसरत।}
\VUsaut{2.4mm}
\VUcentre{10.2pt}{0}{\textit{(एगो लइका के कहल कहानी।)}}
\VUsaut{3.0mm}

%%K t02-18-1 p
१८. --- लचकदार, फुरतीला आ तंदुरुस्त रहे खातिर हमनी के आपन गोड़ आ आपन
बाँह के कसरत करावे के चाहीं। एहीसे हमनी के खेलत बानी जा आ
\VUgras{कसरत} करत बानी जा।

%%K t02-19-1 p
१९. --- हफ्ता भर ई दुनो कसरत इसकूल के अँगना में होला (नकसा नं. १
देखीं)। बाकिर बियफे आ इतवार के हमनी के बड़का घास के मैदान पर जाइले जा,
जहाँ दउड़े आ मन बहलावे खातिर जगहा रहेला।

%%K t02-20-1 p
२०. --- हमनी के मास्टर लोग आ घर के लोग कबो-कबो हमनी के साथे ओहिजा
जालें ; बाकिर कसरत के देखरेख \VUgras{कसरत के मास्टर}
\textsuperscript{(12)} ही करेलें।

%%K t02-21-1 p
२१. --- मैदान में, गेट के बाईं ओर, \VUgras{कसरत के सामान}
\textsuperscript{(1)} बा ; हमनी के \VUgras{सीढ़ी} \textsuperscript{(2)}
से चढ़त बानी जा, \VUgras{छल्ला} \textsuperscript{(3)} आ
\VUgras{ट्रैपीज} \textsuperscript{(4)} पर उलटा लटकत बानी जा, आ
\VUgras{रस्सी के सीढ़ी} \textsuperscript{(5)} भा \VUgras{गाँठ वाला
रस्सी} \textsuperscript{(6)} के चोटी तक हाथ से चढ़ जाइले जा।

%%K t02-22-1 p
२२. --- ओही बेरा हमनी के संगी \VUgras{लकड़ी के घोड़ा}
\textsuperscript{(7)} भा \VUgras{समानांतर डंडा} \textsuperscript{(11)} के
ऊपर से कूदेलें। दोसर लोग \VUgras{मुक्केबाजी} \textsuperscript{(8)} करेला
भा मुखौटा पहिन के \VUgras{पातर तलवार} से \VUgras{तलवारबाजी}
\textsuperscript{(9)} करेला। आ कुछ लोग \VUgras{डंबल उठावेला}
\textsuperscript{(10)}।

%%K t02-tit-5 sub
\VUsaut{4.6mm}
\VUpk{10.2pt}{40}{खेल।}
\VUsaut{3.0mm}

%%K t02-23-1 p
२३. --- कसरत करे वाला लोग के चारो ओर छोटका लइका आ लइकी तरह-तरह के
\VUgras{खेल} खेलेलें। लइकी लोग आपन \VUgras{घेरा} \textsuperscript{(14)}
से मन बहलावेली ; ऊ लोग \VUgras{रस्सी} \textsuperscript{(13)} कूदेली भा
आपन भाई आ भतीजा लोग के \VUgras{क्रोके} \textsuperscript{(15)} के खेल में
बोलावेली। ओह लोग के मजा एहमें आवेला कि जब सामने वाला सोचत रहेला कि ऊ
आसानी से मेहराब पार कर लीही, ओही बेरा ओकर \VUgras{गेंद} के आपन
\VUgras{लकड़ी के हथौड़ी} से ठेल के दूर कर दीहें !

%%K t02-24-1 p
२४. --- छोटका लइकन के ताकत वाला खेल जादे नीक लागेला। ऊ लोग खुसी से
\VUgras{गुल्ली} \textsuperscript{(16)} खेलेला, जेकरा के ऊ लोग
\VUgras{गेंद} \textsuperscript{(17)} से चतुराई से गिरा देला। ऊ लोग
\VUgras{लट्टू} \textsuperscript{(18)} आ \VUgras{बिलबोके}
\textsuperscript{(19)} भा \VUgras{मेढ़क-खेल} \textsuperscript{(20)} खेले
से भी परहेज ना करे। बाकिर ओह लोग के सबसे पसंदीदा खेल \VUgras{गोली}
\textsuperscript{(21)} आ \VUgras{भीत-गेंद} \textsuperscript{(22)} ह,
काहेकि \VUgras{काँच-घर} \textsuperscript{(26)} के भीत हलुक गेंद के वापस
उछाले खातिर बहुते ठीक बा।

%%K t02-25-1 p
२५. --- छोटका योहानेस, आपन \VUgras{बल्ली-गोड़} \textsuperscript{(24)} पर
चलत, बहुते चतुर बा आ बहुते नीक से संतुलन साधेला। ओकरा लगे कुछ संगी ओह
काँच-घर में आ \VUgras{बाड़} के पाछे \VUgras{लुका-छिपी}
\textsuperscript{(25)} खेलत बाड़ें। दोसर लोग \VUgras{थपकी-बूझ}
\textsuperscript{(28)} भा \VUgras{उड़त गेंद} \textsuperscript{(29)} पसंद
करेला, जवन एगो \VUgras{रैकेट} से दोसरा पर उड़त रहेला।

%%K t02-noto-1 noto
\VUnotes{143.2mm}{%
(*) लइकन के खेल के नाँव अकसर पूरा-पूरा देस-विदेस के आपन होला। हमनी के
कोसिस कइल जा कि ओकनी के इडो में जेतना नीक हो सके ओतना नाँव मिले, कबो
ओह काम से जवना से ऊ खेल पहिचानल जाला, त कबो कवनो देस के नाँव लेके, जब
ऊ बहुते बेजोड़ ना होखे। जइसे : \VUgras{पीपा-खेल} (जर्मन आ फ्रेंच) ऊहे
\VUgras{मेढ़क-खेल} \textsuperscript{(20)} (अंगरेजी) ह, जवना में खेले वाला
लोग आपन गोल चकती के छेद वाला पीपा पर ठाढ़ मुँह खोलले लोहा के मेढ़क के
मुँह में फेंके के कोसिस करेला।
\VUgras{कान्ह-कूद} \textsuperscript{(30)} \VUgras{पीठ-कूद} से खाली एतने
में अलग बा : कपार के ऊपर से कूदे खातिर खेले वाला आपन हाथ साथी के कान्ह
पर टेकेला, जबकि « \VUgras{पीठ-कूद} » में ऊ हाथ खाली ओकर पीठ पर टेकेला।
--- भा फेरु कुछ लोग झुकल रहेला आ बाकी लोग ओह लोग के पीठ पर से हाथ
कान्ह पर टेक के कूदेला ; पहिलका लोग के धक्का बिना झुकले सहे के परेला,
दुसरका लोग के संतुलन बिगड़ले बिना कूदे के परेला (नकसा खुदे देखीं)।%
}

%%K t02-26-1 p
२६. --- मैदान के बीच में दू गो गाछ बा। ओहमें से एगो के जड़ लगे सात भा
आठ गो लइका \VUgras{कान्ह-कूद} \textsuperscript{(30)} (*) के खेल जमवले
बाड़ें। ई खेल सब देस में ना जानल जाव ; बाकिर \VUgras{पीठ-कूद} का ह ई सब
जानेला, जवन ऊ लोग एह बेर नइखे खेलत।

%%K t02-tit-6 sub
\VUsaut{5.0mm}
\VUpk{10.2pt}{40}{खुला हवा के खेल।}
\VUsaut{3.4mm}

%%K t02-27-1 p
२७. \VUgras{आँखमिचौली} \textsuperscript{(31)} भी बहुते मजेदार बा ;
लइकी आ लइका बारी-बारी से आपन आँख पर \VUgras{पट्टी} बाँधेलें आ जे
अनाड़ी लोग पकड़ा जाला ओकरा के धर लेलें।

%%K t02-28-1 p
२८. --- मैदान के बीच में एगो बहुते फ्रेंच बाकिर थोड़ा जोरदार खेल बा :
ऊ ह \VUgras{भालू-खेल} \textsuperscript{(32)}, जवन \VUgras{पतंग}
\textsuperscript{(33)} के ओतना सांत खेल से भा अंगरेजी खेल
\VUgras{टेनिस} \textsuperscript{(34)} से भी बहुते जादे सोर करेला।

%%K t02-29-1 p
२९. --- ई आखिरी खेल बड़का लइकन के खुसी ह। हमनी के मास्टर लोग खुदे एकरा
के खुसी से खेलेला। एहमें बहुते चतुराई आ फुरती चाहीं आ हमरा ई बहुते नीक
लागेला। \VUgras{झूला} \textsuperscript{(35)} के खेल हम लइकी लोग खातिर
छोड़ देत बानी।

%%K t02-30-1 p
३०. --- हमरा एगो दोसर अंगरेजी खेल नीक ना लागे, जवन सायद खतरनाक हो
जाव।

%%K t02-30-1b p
हम \VUgras{फुटबॉल} \textsuperscript{(36)} के बारे में कहे चाहत बानी,
जवन हमार बड़का भाई के खुसी देला। हमरा त आपन पुरनका \VUgras{बार के खेल}
\textsuperscript{(38)} जादे नीक लागेला, जवन ओतने सेहतमंद बा आ सचहूँ कम
जंगली बा।

%%K t02-tit-7 sub
\VUsaut{4.6mm}
\VUpk{10.2pt}{40}{साइकिल आ गेंद।}
\VUsaut{3.0mm}

%%K t02-31-1 p
३१. --- हम \VUgras{साइकिल} \textsuperscript{(37)} से मैदान के चक्कर
लगावल भी खुसी से करत बानी।

%%K t02-32-1 p
३२. --- अगिला साल हम \VUgras{घुड़सवारी} \textsuperscript{(39)} सीखब।
कतना सुंदर लागेला ऊ घोड़ा जवन ठसक से चलेला भा दुलकी भरेला, आ जवन सरपट
दउड़त बाधा के ऊपर से कूद जाला !

%%K t02-33-1 p
३३. --- गरमी में हमनी के \VUgras{तैराकी} \textsuperscript{(40)} सीखत
बानी जा। हमनी के मजा लेके नदी के साफ आ ठंढा पानी में डुबकी लगावत आ
नहात बानी जा। कबो-कबो हमनी के कागज के \VUgras{गुब्बारा}
\textsuperscript{(41)} भी उड़ाइले जा, जवन गरम हवा से भरते शान से आसमान
में चढ़ जाला।

%%K t02-34-1 p
३४. --- आउर भी बहुते खेल बा, बाकिर ओकनी के गिनल हमरा से खतम ना होई, आ
एह सार से पता चल जाला कि बियफे के हमनी के सुंदर मैदान पर मन ना ऊबे।

%%K t02-34-2 p
टीप। --- \textit{मास्टर साहेब एह अध्याय के आसानी से पूरा कर लीहें ;
हर खास खेल के आउर बिस्तार से बयान करके।}
